% from ega2-1-fr.tex
\oldpage[II]{27}
qui précède prouve donc, compte tenu de l'axiome $(T_0)$ valable dans
$D_+(f)$ et dans $\operatorname{Spec}(S_{(f)})$, que $\psi_f$ est injective
et que l'application réciproque
$\psi_f(D_+(f))\longrightarrow D_+(f)$ est continue. Enfin, pour voir que
$\psi_f$ est surjective, on remarque que, si $\mathfrak{q}_0$ est un idéal
premier de $S_{(f)}$, et si, pour tout $n>0$, on désigne par
$\mathfrak{p}_n$ l'ensemble des $x\in S_n$ tels que
$x^d/f^n\in\mathfrak{q}_0$, les $\mathfrak{p}_n$ vérifient les conditions
de (\hyperref[II.2.1.9-fr]{2.1.9})~: en effet, si $x\in S_n$, $y\in S_n$
sont tels que $x^d/f^n\in\mathfrak{q}_0$ et
$y^d/f^n\in\mathfrak{q}_0$, on a
$(x+y)^{2d}/f^{2n}\in\mathfrak{q}_0$, d'où
$(x+y)^d/f^n\in\mathfrak{q}_0$ puisque $\mathfrak{q}_0$ est premier~;
cela prouve que les $\mathfrak{p}_n$ sont des sous-groupes des $S_n$, et la
vérification des autres conditions de
(\hyperref[II.2.1.9-fr]{2.1.9}) est immédiate, compte tenu de ce que
$\mathfrak{q}_0$ est premier. Si $\mathfrak{p}$ est l'idéal premier gradué
de $S$ ainsi défini, on a bien
$\psi_f(\mathfrak{p})=\mathfrak{q}_0$, car si $x\in S_{nd}$, les relations
$x/f^n\in\mathfrak{q}_0$ et $x^d/f^{nd}\in\mathfrak{q}_0$ sont
équivalentes, $\mathfrak{q}_0$ étant premier.
\end{proof}

\begin{corollary}[2.3.7]
\label{II.2.3.7-fr}
Pour que $D_+(f)=\varnothing$, il faut et il suffit que $f$ soit nilpotent.
\end{corollary}

\begin{proof}
En effet, pour que $\operatorname{Spec}(S_{(f)})=\varnothing$, il faut et
il suffit que $S_{(f)}=0$, ou encore que $1=0$ dans $S_f$, ce qui signifie
par définition que $f$ est nilpotent.
\end{proof}

\begin{corollary}[2.3.8]
\label{II.2.3.8-fr}
Soit $E$ une partie de $S_+$. Les conditions suivantes sont équivalentes~:
\begin{enumerate}
  \item[\rm a)] $V_+(E)=X=\operatorname{Proj}(S)$.
  \item[\rm b)] Tout élément de $E$ est nilpotent.
  \item[\rm c)] Les composants homogènes de tout élément de $E$ sont
    nilpotents.
\end{enumerate}
\end{corollary}

\begin{proof}
Il est clair que c) entraîne b) et que b) entraîne a). Si $\mathfrak{J}$ est
l'idéal gradué de $S$ engendré par $E$, la condition a) équivaut à
$V_+(\mathfrak{J})=X$~; \emph{a fortiori}, a) entraîne que tout élément
homogène $f\in\mathfrak{J}$ est tel que $V_+(f)=X$, donc $f$ est nilpotent
en vertu de (\hyperref[II.2.3.7-fr]{2.3.7}).
\end{proof}

\begin{corollary}[2.3.9]
\label{II.2.3.9-fr}
Si $\mathfrak{J}$ est un idéal gradué de $S_+$,
$r_+(\mathfrak{J})$ est l'intersection des idéaux premiers gradués de
$S_+$ qui contiennent $\mathfrak{J}$.
\end{corollary}

\begin{proof}
En considérant l'anneau gradué $S/\mathfrak{J}$, on peut se ramener au cas
où $\mathfrak{J}=0$. Il faut prouver que si $f\in S_+$ n'est pas
nilpotent, il y a un idéal premier gradué de $S$ ne contenant pas $f$~; or,
un au moins des composants homogènes de $f$ n'est pas nilpotent, et on peut
donc supposer $f$ homogène~; l'assertion résulte alors de
(\hyperref[II.2.3.7-fr]{2.3.7}).
\end{proof}

\begin{env}[2.3.10]
\label{II.2.3.10-fr}
Pour toute partie $Y$ de $X=\operatorname{Proj}(S)$, soit
$\mathfrak{j}_+(Y)$ l'ensemble des $f\in S_+$ tels que
$Y\subset V_+(f)$~; il revient au même de dire que
$\mathfrak{j}_+(Y)=\mathfrak{j}(Y)\cap S_+$~; $\mathfrak{j}_+(Y)$ est
donc un idéal de $S_+$, égal à sa racine dans $S_+$.
\end{env}

\begin{proposition}[2.3.11]
\label{II.2.3.11-fr}
\begin{enumerate}
  \item[\rm(i)] Pour toute partie $E$ de $S_+$,
    $\mathfrak{j}_+(V_+(E))$ est la racine dans $S_+$ de l'idéal gradué de
    $S_+$ engendré par $E$.
  \item[\rm(ii)] Pour toute partie $Y$ de $X$,
    $V_+(\mathfrak{j}_+(Y))=\overline{Y}$, adhérence de $Y$ dans $X$.
\end{enumerate}
\end{proposition}

\begin{proof}
\begin{enumerate}
  \item[\rm(i)] Si $\mathfrak{J}$ est l'idéal gradué de $S_+$ engendré
    par $E$, on a $V_+(E)=V_+(\mathfrak{J})$ et l'assertion résulte alors de
    (\hyperref[II.2.3.9-fr]{2.3.9}).
  \item[\rm(ii)] Comme
    $V_+(\mathfrak{J})=\bigcap_{f\in\mathfrak{J}}V_+(f)$, la relation
    $Y\subset V_+(\mathfrak{J})$ entraîne $Y\subset V_+(f)$ pour tout
    $f\in\mathfrak{J}$, et par suite
    $\mathfrak{j}_+(Y)\supset\mathfrak{J}$, d'où
    $V_+(\mathfrak{j}_+(Y))\subset V_+(\mathfrak{J})$, ce qui prouve (ii)
    par définition des ensembles fermés.
\end{enumerate}
\end{proof}

\begin{corollary}[2.3.12]
\label{II.2.3.12-fr}
Les parties fermées $Y$ de $X=\operatorname{Proj}(S)$ et les idéaux gradués
de $S_+$ égaux à leur racine dans $S_+$ se correspondent biunivoquement par
les applications décroissantes
$Y\longmapsto\mathfrak{j}_+(Y)$, $\mathfrak{J}\longmapsto
V_+(\mathfrak{J})$~; à la réunion $Y_1\cup Y_2$ de deux parties fermées de
$X$ correspond
